\chapter{コンピュータとは}
\label{chap:computer}

% この章で学ぶこと
\begin{goalbox}
  \begin{itemize}
    \item コンピュータを構成する主な部品（CPU・メモリ・ストレージ・GPU・入出力装置）の役割
    \item それぞれの部品が連携してプログラムを実行するしくみ
    \item ハードウェアとソフトウェア、プログラムとアルゴリズム、そして OS の役割
    \item ロボットに載っているコンピュータの種類
  \end{itemize}
\end{goalbox}

ロボットは、カメラやセンサで周りの様子を読み取り、どう動くかを考え、モータを動かします。
この「考える」部分を担っているのがコンピュータです。
本書ではこれから Linux や ROS 2 を使ってロボットを動かしていきますが、その前に、コンピュータがどのような部品でできていて、どのように計算しているのかを簡単に見ておきましょう。
ここで紹介する内容は、後の章で出てくる「CPU 使用率」「メモリ不足」「GPU が必要」といった話を理解する土台になります。

\section{コンピュータの構成要素}
\label{sec:computer-parts}

ノートパソコンもデスクトップパソコンも、ロボットに載っている小さなコンピュータも、基本的な構成は同じです。
コンピュータは主に、次の部品でできています（図\ref{fig:computer-parts}）。

\begin{figure}[tb]
  \centering
  % 図：コンピュータの主な構成要素（chapters/linux/computer.tex）
\begin{tikzpicture}
  % バス
  \draw[line width=4pt, black!25] (-0.4,0) -- (9.4,0);
  \node[figlabel, anchor=north west, text=black!60] at (2.3,-0.15) {バス（マザーボード上の配線）};
  % 上段
  \node[figpart, minimum width=24mm] (cpu) at (1.4,1.5) {\textbf{CPU}\\{\footnotesize 命令の実行}};
  \node[figpart, minimum width=24mm] (mem) at (4.5,1.5) {\textbf{メモリ}\\{\footnotesize 作業中のデータ}};
  \node[figpart, minimum width=24mm] (gpu) at (7.6,1.5) {\textbf{GPU}\\{\footnotesize 大量の並列計算}};
  % 下段
  \node[figpart, minimum width=24mm] (sto) at (1.4,-1.7) {\textbf{ストレージ}\\{\footnotesize データの保存}};
  \node[figpart, minimum width=40mm] (io) at (7.0,-1.7) {\textbf{入出力装置}\\{\footnotesize キーボード・ディスプレイ}\\{\footnotesize カメラ・センサ・モータ}};
  % 接続
  \foreach \n in {cpu,mem,gpu} \draw[{Stealth[length=1.6mm]}-{Stealth[length=1.6mm]}, thick] (\n.south) -- (\n.south |- 0,0.08);
  \foreach \n in {sto,io} \draw[{Stealth[length=1.6mm]}-{Stealth[length=1.6mm]}, thick] (\n.north) -- (\n.north |- 0,-0.08);
\end{tikzpicture}

  \caption{コンピュータの主な構成要素}
  \label{fig:computer-parts}
\end{figure}

ここでは、それぞれの部品の役割を、机で作業をする人にたとえて説明します。

\subsection{CPU：計算と命令の実行を担う頭脳}

\textbf{CPU}（Central Processing Unit、中央処理装置）は、プログラムに書かれた命令を 1 つずつ読み取り、計算や判断を行う部品です。
コンピュータのすべての処理は、最終的に CPU が実行します。
机で作業をする人にたとえると、CPU は「作業をする人そのもの」です。

CPU の性能は、主に次の 2 つで表されます。

\begin{description}
  \item[クロック周波数]
    1 秒間に何回処理のタイミングを刻むかを表します。
    単位は GHz（ギガヘルツ）で、3 GHz なら 1 秒間に 30 億回です。
    数字が大きいほど、1 つ 1 つの処理を速くこなせます。
  \item[コア数]
    CPU の中で、独立して処理を実行できる部分を\textbf{コア}といいます。
    最近の CPU は、1 つの CPU の中に数個から数十個のコアを持っています。
    コアが多いほど、たくさんの処理を同時にこなせます。
\end{description}

CPU は、複雑な判断を含む処理を、順番に素早くこなすことが得意です。

\subsection{メモリ：作業中のデータを置く机}

\textbf{メモリ}（主記憶装置、RAM とも呼ばれます）は、CPU が今まさに使っているプログラムやデータを一時的に置いておく場所です。
CPU はメモリに置かれたデータを読み書きしながら処理を進めます。
たとえると、メモリは「作業をするための机」です。

メモリには、次のような特徴があります。

\begin{itemize}
  \item 読み書きがとても速い
  \item 電源を切ると、中身が消えてしまう
  \item 容量は、パソコンでは 8〜32 GB（ギガバイト）程度が一般的
\end{itemize}

机が広いほどたくさんの資料を広げて作業できるように、メモリの容量が大きいほど、たくさんのプログラムを同時に動かしたり、大きなデータを扱ったりできます。
メモリが足りなくなると、コンピュータの動作が極端に遅くなったり、プログラムが強制的に終了させられたりします。

\subsection{ストレージ：データをしまっておく本棚}

\textbf{ストレージ}（補助記憶装置）は、プログラムやデータを長期間保存しておく場所です。
SSD（Solid State Drive）や HDD（Hard Disk Drive）がこれにあたります。
たとえると、ストレージは「資料をしまっておく本棚」です。

ストレージには、次のような特徴があります。

\begin{itemize}
  \item 電源を切っても、中身が消えない
  \item 容量が大きい（パソコンでは数百 GB〜数 TB（テラバイト）程度）
  \item メモリに比べると、読み書きがずっと遅い
\end{itemize}

OS やアプリケーション、作成したファイルは、すべてストレージに保存されています。

\subsection{GPU：大量の計算を同時にこなす専門家}

\textbf{GPU}（Graphics Processing Unit）は、もともとは画面に表示する画像を作るための部品です。
3D のゲームの画面のように、たくさんの点（画素）の色を一斉に計算する処理を得意としています。

GPU は、CPU に比べて 1 つ 1 つのコアの能力は低いものの、数千個ものコアを持っています。
そのため、「同じ種類の単純な計算を、大量のデータに対して一斉に行う」処理を CPU よりもはるかに速くこなせます（図\ref{fig:computer-cpu-gpu}）。

\begin{figure}[tb]
  \centering
  % 図：CPU と GPU の違い（chapters/linux/computer.tex）
\begin{tikzpicture}
  % CPU：少数の高性能なコア
  \node[figbox, minimum width=38mm, minimum height=30mm, fill=blue!3] (cpu) at (0,0) {};
  \node[font=\small\bfseries, anchor=north] at (cpu.north) {CPU};
  \foreach \x in {0,1} \foreach \y in {0,1}
    \node[draw, fill=blue!20, minimum size=9mm, font=\scriptsize] at (-0.55+1.1*\x, -0.85+1.0*\y) {コア};
  \node[figlabel, below=2mm of cpu] {少数の高性能なコア\\複雑な処理を順番にこなす};
  % GPU：多数の単純なコア
  \node[figbox, minimum width=50mm, minimum height=30mm, fill=orange!3] (gpu) at (5.2,0) {};
  \node[font=\small\bfseries, anchor=north] at (gpu.north) {GPU};
  \foreach \x in {0,...,15} \foreach \y in {0,...,6}
    \fill[orange!45] (3.25+0.25*\x, -1.3+0.27*\y) rectangle ++(0.19,0.2);
  \node[figlabel, below=2mm of gpu] {多数の単純なコア\\単純な計算を一斉にこなす};
\end{tikzpicture}

  \caption{CPU と GPU の違い}
  \label{fig:computer-cpu-gpu}
\end{figure}

たとえると、CPU が「どんな仕事でもこなせる数人の熟練した職人」だとすれば、GPU は「単純な作業を一斉にこなす大勢の作業員」です。
この特徴から、GPU は現在では画像の表示だけでなく、次のような用途にも使われています。

\begin{itemize}
  \item カメラ画像の処理や、深層学習（ディープラーニング）による物体認識
  \item ロボットや周囲の環境の 3D 表示（\ref{chap:tf_viz}章で扱う RViz）
  \item 物理シミュレーション（\ref{chap:simulation}章で扱う Gazebo）
\end{itemize}

GPU は、CPU とは別に専用のメモリ（\textbf{VRAM}、ビデオメモリ）を持っているのが一般的です。
GPU で計算するときは、データをメモリから VRAM に送り、計算が終わったら結果をメモリに戻します。

\subsection{入出力装置：外の世界とつながる窓口}

コンピュータに情報を入力したり、コンピュータから情報を出力したりする装置を\textbf{入出力装置}といいます。
パソコンでは、キーボードやマウスが入力装置、ディスプレイやスピーカが出力装置です。

ロボットでは、カメラや LiDAR（レーザで周囲の距離を測るセンサ）、IMU（加速度や角速度を測るセンサ）などのセンサが入力装置、モータやスピーカが出力装置にあたります。
これらの装置は、USB やネットワーク（LAN）などを通じてコンピュータにつながっています。

\begin{tablebox}[コンピュータの主な部品とその役割][tab:computer-parts]
  \begin{tabular}{llp{0.45\linewidth}}
    \toprule
    部品 & たとえ & 役割 \\
    \midrule
    CPU & 作業をする人 & 命令を実行し、計算や判断をする \\
    メモリ & 机 & 作業中のデータを一時的に置く \\
    ストレージ & 本棚 & データを長期間保存する \\
    GPU & 大勢の作業員 & 単純で大量の計算を一斉にこなす \\
    入出力装置 & 目・耳・手足 & 外の世界と情報をやり取りする \\
    \bottomrule
  \end{tabular}
\end{tablebox}

\section{各部品が連携して動くしくみ}
\label{sec:computer-flow}

それでは、これらの部品がどのように連携してプログラムを実行するのかを見てみましょう。
ここでは例として、カメラの画像から人を見つけるプログラムを実行する場面を考えます（図\ref{fig:computer-flow}）。

\begin{figure}[tb]
  \centering
  % 図：プログラムが実行されるときのデータの流れ（chapters/linux/computer.tex）
\begin{tikzpicture}[node distance=1.2cm]
  \node[figpart, minimum width=22mm] (mem) at (3.7,0) {\textbf{メモリ}};
  \node[figpart, minimum width=22mm] (cpu) at (3.7,2.2) {\textbf{CPU}};
  \node[figpart, minimum width=22mm] (sto) at (-0.4,0) {\textbf{ストレージ}};
  \node[figpart, minimum width=22mm] (gpu) at (7.8,0) {\textbf{GPU}\\{\footnotesize（VRAM）}};
  \node[figbox, minimum width=22mm] (cam) at (3.7,-2.2) {カメラ};
  \node[figbox, minimum width=22mm] (out) at (7.8,2.2) {ディスプレイ\\モータ};
  % (1) 読み込み / (6) 保存
  \draw[figarrow] ([yshift=2mm]sto.east) -- node[figlabel, above]{\tikz\node[fignum]{1}; 読み込み} ([yshift=2mm]mem.west);
  \draw[figarrow] ([yshift=-2mm]mem.west) -- node[figlabel, below]{\tikz\node[fignum]{6}; 保存} ([yshift=-2mm]sto.east);
  % (2) 命令の実行
  \draw[figarrow, {Stealth[length=2mm]}-{Stealth[length=2mm]}] (mem.north) -- node[figlabel, right]{\tikz\node[fignum]{2}; 命令・データ} (cpu.south);
  % (3) 画像の入力
  \draw[figarrow] (cam.north) -- node[figlabel, left]{\tikz\node[fignum]{3}; 画像} (mem.south);
  % (4) GPU に計算を任せる
  \draw[figarrow] ([yshift=2mm]mem.east) -- node[figlabel, above]{\tikz\node[fignum]{4}; 計算} ([yshift=2mm]gpu.west);
  \draw[figarrow] ([yshift=-2mm]gpu.west) -- node[figlabel, below]{結果} ([yshift=-2mm]mem.east);
  % (5) 出力
  \draw[figarrow] (cpu.east) -- node[figlabel, above]{\tikz\node[fignum]{5}; 出力} (out.west);
\end{tikzpicture}

  \caption{プログラムが実行されるときのデータの流れ}
  \label{fig:computer-flow}
\end{figure}

\begin{enumerate}
  \item \textbf{プログラムを読み込む}：
    プログラムはふだんストレージに保存されています。
    実行するとき、まずストレージからメモリに読み込まれます（本棚から資料を取り出して、机に広げる）。
  \item \textbf{CPU が命令を実行する}：
    CPU は、メモリに置かれたプログラムの命令を 1 つずつ読み取って実行します。
  \item \textbf{データを受け取る}：
    カメラで撮影した画像のデータが、USB などを通じてメモリに届きます。
  \item \textbf{重い計算を GPU に任せる}：
    画像から人を見つける処理には、膨大な量の計算が必要です。
    CPU は画像のデータを GPU に送り、計算を任せます。
    GPU は多数のコアで一斉に計算し、結果を返します。
  \item \textbf{結果を使って判断し、出力する}：
    CPU は GPU から返ってきた結果をもとに「人がいるので止まる」といった判断をし、ディスプレイに結果を表示したり、モータに指令を送ったりします。
  \item \textbf{データを保存する}：
    記録しておきたいデータは、メモリからストレージに書き込みます（机の上の資料を本棚にしまう）。
\end{enumerate}

このように、CPU が全体の司令塔となって処理を進め、メモリを作業場所として使い、必要に応じて GPU に計算を任せたり、ストレージや入出力装置とデータをやり取りしたりしています。

\begin{point}[速さと容量の関係]
  データを置く場所は、速いものほど容量が小さく、遅いものほど容量が大きくなっています。
  \begin{itemize}
    \item メモリ：速いが容量は小さい（数 GB〜数十 GB）。電源を切ると消える
    \item ストレージ：遅いが容量は大きい（数百 GB〜数 TB）。電源を切っても残る
  \end{itemize}
  そのため、よく使うデータはメモリに、長く保存したいデータはストレージに置く、という使い分けをしています。
  \ref{chap:process}章では、コマンドを使って CPU やメモリがどれくらい使われているのかを確認する方法を学びます。
\end{point}

\section{ハードウェアとソフトウェア}
\label{sec:computer-software}

ここまで紹介した CPU やメモリのような、手で触れることのできる物理的な部品を\textbf{ハードウェア}といいます。
これに対して、コンピュータに何をさせるのかを指示するプログラムのことを\textbf{ソフトウェア}といいます。
ハードウェアだけでは何もできず、ソフトウェアがあって初めてコンピュータは役に立つ仕事をします。

\subsection{プログラムとは}

\textbf{プログラム}は、コンピュータに実行させる命令を順番に並べたものです。
コンピュータは、とても速く、正確に計算できますが、自分で考えて動くことはできません。
「何を」「どの順番で」「どういうときに」行うのかを、すべて人間が指示してあげる必要があります。
その指示書がプログラムで、プログラムを書くことを\textbf{プログラミング}といいます。

料理にたとえると、プログラムはレシピのようなものです。
ただし、コンピュータは、人間のように「適量」「いい感じに」といったあいまいな指示を理解できません。
「塩を 2 g 入れる」「中火で 3 分焼く」のように、誰が実行しても同じ結果になるくらい、細かく正確に書く必要があります。

この本でこれから使っていく \cmd{ls} や \cmd{git} などのコマンドも、ROS 2 も、誰かが書いたプログラムです。
そしてロボットは、プログラムに書かれた指示に従って、センサの値を読み、行動を決め、モータを動かしています。
ロボットを思いどおりに動かせるようになるには、自分でプログラムを書けるようになる必要があるのです。

CPU が直接理解できるのは、0 と 1 の並びで表された\textbf{機械語}だけです。
しかし、人間が機械語でプログラムを書くのはとても大変です。
そこで人間は、Python や C++ のような\textbf{プログラミング言語}を使って、人間にも読みやすい形でプログラムを書きます。
書いたプログラムは、専用のソフトウェアによって機械語に翻訳されてから、CPU で実行されます（\ref{sec:python-compile-interpret}節）。
本書の後半では、Python を使って ROS 2 のプログラムを書きます（\ref{chap:rclpy}章）。

\subsection{アルゴリズム}

ある問題を解くための手順のことを\textbf{アルゴリズム}といいます。
たとえば、「障害物にぶつからないようにロボットを動かす」という問題に対しては、次のような手順が考えられます。

\begin{enumerate}
  \item 距離センサで、前にある障害物までの距離を測る
  \item 距離が 0.3 m より小さければモータを止め、そうでなければ前に進む
  \item 1 に戻って、これを繰り返す
\end{enumerate}

プログラミングでは、まずこのように「どういう手順で解くか」（アルゴリズム）を考え、それをプログラミング言語で書き表します。
いきなりプログラムを書き始めるのではなく、手順を日本語や図で整理してから書くと、うまくいきやすくなります。

\subsection{順次・分岐・繰り返し}

上の手順を、図にしてみましょう（図\ref{fig:computer-flowchart}）。
このように、処理の流れを図で表したものを\textbf{流れ図}（フローチャート）といいます。

\begin{figure}[tb]
  \centering
  % 図：障害物の前で止まるロボットのプログラムの流れ図（chapters/linux/computer.tex）
\begin{tikzpicture}[
    node distance=7mm,
    proc/.style={figbox, minimum width=34mm},
    term/.style={figbox, rounded corners=3mm, minimum width=24mm, fill=black!6},
    decision/.style={draw, diamond, aspect=2.4, align=center, font=\small, fill=orange!10, inner sep=1pt}]
  \node[term] (start) {開始};
  \node[proc, below=of start] (read) {距離センサの値を読む};
  \node[decision, below=of read] (check) {距離が 0.3 m\\より小さい？};
  \node[proc, below left=9mm and -4mm of check] (stop) {モータを止める};
  \node[proc, below right=9mm and -4mm of check] (go) {前に進む};
  \draw[figarrow] (start) -- (read);
  \draw[figarrow] (read) -- (check);
  \draw[figarrow] (check.west) -| node[figlabel, above, pos=0.25]{はい} (stop.north);
  \draw[figarrow] (check.east) -| node[figlabel, above, pos=0.25]{いいえ} (go.north);
  \coordinate (join) at ($(stop.south)!0.5!(go.south) + (0,-0.6)$);
  \draw[thick] (stop.south) |- (join);
  \draw[thick] (go.south) |- (join);
  \draw[figarrow] (join) -- ++(0,-0.3) -- ++(-4.2,0) |- node[figlabel, left, pos=0.25, text=green!50!black]{繰り返し} (read.west);
  \node[figlabel, text=blue!60!black, anchor=west] at ([xshift=2mm]read.east) {順次};
  \node[figlabel, text=orange!70!black, anchor=west] at ([xshift=1mm]check.east |- check.north) {分岐};
\end{tikzpicture}

  \caption{障害物の前で止まるロボットのプログラムの流れ図}
  \label{fig:computer-flowchart}
\end{figure}

この図には、プログラムの基本となる 3 つの流れがすべて含まれています。

\begin{description}
  \item[順次] 上から順番に、1 つずつ処理を実行する（「センサの値を読む」→「判断する」）
  \item[分岐] 条件によって、実行する処理を選ぶ（距離によって「止める」か「進む」かを選ぶ）
  \item[繰り返し] 同じ処理を何度も実行する（センサを読むところに戻って、ずっと繰り返す）
\end{description}

実は、どんなに複雑なプログラムも、この\textbf{順次・分岐・繰り返し}の 3 つを組み合わせて作られています。
ロボットのプログラムも、ゲームも、OS も同じです。
本書では、シェルスクリプト（\ref{chap:shellscript}章）や Python（\ref{chap:python_basics}章）で、分岐や繰り返しの書き方を学びます。
書き方は言語によって違っても、プログラムの考え方は共通です。

\subsection{バグとデバッグ}

プログラムは、書いたとおりにしか動きません。
指示が間違っていたり、考えていなかった場合があったりすると、プログラムは思ったとおりに動きません。
このようなプログラムの誤りを\textbf{バグ}といい、バグを見つけて直すことを\textbf{デバッグ}といいます（\ref{sec:python-class-debug}節）。

最初から完璧なプログラムを書ける人はいません。
プログラミングは、「書いて、動かして、間違いを見つけて、直す」の繰り返しです。
エラーが出ても落ち込まずに、少しずつ動くプログラムに近づけていきましょう。

\subsection{ソフトウェアの種類}

ソフトウェアは、大きく次の 2 つに分けられます。

\begin{description}
  \item[アプリケーションソフトウェア]
    Web ブラウザや文書作成ソフト、ゲームのように、利用者が目的を果たすために使うソフトウェアです。
    ロボットを動かすプログラムも、アプリケーションソフトウェアの一種です。
  \item[基本ソフトウェア（OS）]
    ハードウェアを管理し、アプリケーションソフトウェアが動く土台となるソフトウェアです。
    次の節で詳しく説明します。
\end{description}

\section{OS とは何か}
\label{sec:computer-os}

\textbf{OS}（Operating System）は、ハードウェアとアプリケーションソフトウェアの間に立って、コンピュータ全体を管理するソフトウェアです。
パソコン用の Windows や macOS、スマートフォン用の Android や iOS、そして本書で扱う Linux が OS です。
コンピュータの電源を入れると、まず OS が起動し、その上でさまざまなアプリケーションが動きます（図\ref{fig:computer-os-layer}）。

\begin{figure}[tb]
  \centering
  % 図：ハードウェア・OS・アプリケーションの関係（chapters/linux/computer.tex）
\begin{tikzpicture}[x=1cm, y=1cm]
  \def\w{9.6}
  % ハードウェア
  \node[figbox, minimum width=\w cm, minimum height=16mm, fill=black!6, anchor=south west] (hw) at (0,0) {};
  \node[font=\small\bfseries, anchor=north west] at (hw.north west) {ハードウェア};
  \foreach \t [count=\i from 0] in {CPU,メモリ,ストレージ,GPU,入出力装置}
    \node[figbox, minimum width=17mm, minimum height=6mm, font=\footnotesize] at (1.0+1.9*\i, 0.5) {\t};
  % OS
  \node[figpart, minimum width=\w cm, minimum height=12mm, anchor=south west, fill=blue!12] (os) at (0,1.8)
    {\textbf{OS}（Linux など）\\{\footnotesize CPU・メモリ・ファイル・周辺機器の管理}};
  % アプリケーション
  \node[figbox, minimum width=\w cm, minimum height=16mm, fill=green!6, anchor=south west] (app) at (0,3.2) {};
  \node[font=\small\bfseries, anchor=north west] at (app.north west) {アプリケーション};
  \foreach \t [count=\i from 0] in {Web ブラウザ,エディタ,ROS 2 のプログラム}
    \node[figbox, minimum width=28mm, minimum height=6mm, font=\footnotesize] at (1.7+3.1*\i, 3.7) {\t};
  % 利用者
  \node[figbox, minimum width=30mm] (user) at (\w/2, 6.0) {利用者};
  \draw[figarrow, {Stealth[length=2mm]}-{Stealth[length=2mm]}] (user.south) -- (app.north -| user.south);
\end{tikzpicture}

  \caption{ハードウェア・OS・アプリケーションの関係}
  \label{fig:computer-os-layer}
\end{figure}

\subsection{OS の役割}

OS の主な役割は、次のとおりです。

\begin{description}
  \item[CPU の割り当て]
    コンピュータでは、たくさんのプログラムが同時に動いています。
    OS は、どのプログラムにどれだけ CPU を使わせるかを細かく切り替えながら管理しています。
    これにより、音楽を聴きながら Web ページを見る、といったことができます。
  \item[メモリの管理]
    それぞれのプログラムにメモリの領域を割り当て、ほかのプログラムの領域を勝手に書き換えないように守ります。
  \item[ファイルの管理]
    ストレージに保存されたデータを「ファイル」や「ディレクトリ」という形で整理し、名前で読み書きできるようにします（\ref{chap:files}章）。
  \item[周辺機器の管理]
    キーボード、ディスプレイ、カメラ、USB 機器などの周辺機器を制御します。
    周辺機器を動かすためのソフトウェアを\textbf{デバイスドライバ}といいます。
  \item[利用者とのやり取り]
    利用者がコンピュータを操作するための画面や操作方法を提供します。
    マウスで操作する GUI と、キーボードでコマンドを入力する CLI があります（\ref{chap:terminal_shell}章）。
\end{description}

\subsection{OS があることの利点}

もし OS がなかったら、アプリケーションを作る人は、使うコンピュータのハードウェアに合わせて、CPU やメモリ、周辺機器を直接操作するプログラムを書かなければなりません。
OS がハードウェアの違いを吸収し、アプリケーションに共通の使い方を提供してくれるおかげで、同じアプリケーションをさまざまなコンピュータで動かすことができます。

ROS 2 も、OS の上で動くソフトウェアです。
ROS 2 は主に Linux の上で開発されており、本書でも Linux（Ubuntu）の上で ROS 2 を使います。
次の\ref{chap:what_is_linux}章では、その Linux とはどのような OS なのかを見ていきます。

\section{ロボットに載っているコンピュータ}
\label{sec:computer-robot}

ロボットには、目的や大きさに応じてさまざまなコンピュータが使われます。
大きなロボットでは、役割の違う複数のコンピュータを組み合わせることもよくあります（図\ref{fig:computer-robot}）。

\begin{figure}[tb]
  \centering
  % 図：ロボットに載っているコンピュータの構成例（chapters/linux/computer.tex）
\begin{tikzpicture}
  \node[figbox, minimum width=17mm] (cam) at (0,0.8) {カメラ};
  \node[figbox, minimum width=17mm] (lidar) at (0,-0.8) {LiDAR};
  \node[figpart, minimum width=30mm, minimum height=22mm] (main) at (3.7,0)
    {\textbf{メインの}\\\textbf{コンピュータ}\\{\footnotesize Linux + ROS 2}\\{\footnotesize （GPU 付き）}};
  \node[figpart, minimum width=19mm, fill=orange!10] (mcu) at (7.3,0) {\textbf{マイコン}};
  \node[figbox, minimum width=19mm] (motor) at (7.3,-1.8) {モータ};
  \draw[figarrow] (cam.east) -- node[figlabel, above, sloped]{USB} (main.west |- cam.east);
  \draw[figarrow] (lidar.east) -- node[figlabel, below, sloped]{LAN} (main.west |- lidar.east);
  \draw[figarrow] (main.east) -- node[figlabel, above]{USB}node[figlabel, below]{指令} (mcu.west);
  \draw[figarrow] (mcu.south) -- node[figlabel, right]{制御信号} (motor.north);
  \node[figlabel, below=1mm of main, text=black!70] {センサの処理\\行動の判断};
  \node[figlabel, above=1mm of mcu, text=black!70] {モータの制御};
\end{tikzpicture}

  \caption{ロボットに載っているコンピュータの構成例}
  \label{fig:computer-robot}
\end{figure}

\begin{description}
  \item[パソコン（ノートパソコン・小型パソコン）]
    一般的なパソコンをロボットに載せて使います。
    性能が高く、Linux と ROS 2 をそのまま動かせるので、研究や競技用のロボットでよく使われます。
  \item[シングルボードコンピュータ]
    Raspberry Pi のように、1 枚の小さな基板にコンピュータの機能をまとめたものです。
    小型で消費電力が少なく、Linux を動かすことができます。
  \item[GPU を搭載した組込み向けコンピュータ]
    NVIDIA Jetson のように、小型ながら GPU を搭載したコンピュータです。
    カメラ画像を使った物体認識など、GPU を必要とする処理をロボットの上で行いたいときに使います。
  \item[マイコン（マイクロコントローラ）]
    CPU・メモリ・入出力の機能を 1 つのチップにまとめた、小さく安価なコンピュータです。
    Arduino などが有名です。
    性能はパソコンに比べてずっと低いものの、モータの回転を細かいタイミングで制御するといった、決まった処理を正確な時間で繰り返すことが得意です。
    多くの場合、Linux のような OS は使わずに動かします。
\end{description}

典型的な移動ロボットでは、メインのコンピュータ（パソコンなど）で Linux と ROS 2 を動かしてセンサの処理や行動の判断を行い、マイコンがメインのコンピュータからの指令を受けてモータを制御する、という役割分担がよく見られます。
本書で学ぶのは、主にこのメインのコンピュータの上での開発です。

\begin{column}{パソコンのスペック表の読み方}
  パソコンのカタログには、「CPU：8 コア 3.0 GHz、メモリ：16 GB、ストレージ：SSD 512 GB、GPU：あり」のようなスペック（性能）が書かれています。
  この章で学んだ知識があれば、それぞれが何を意味しているのかがわかるはずです。
  ROS 2 での開発では、シミュレータや 3D 表示を使うため、メモリには余裕があるほうが快適です。
  また、\ref{chap:simulation}章で扱う Gazebo のようなシミュレータを快適に動かしたい場合や、深層学習を使いたい場合は、GPU を搭載したコンピュータがあると便利です。
\end{column}

\section*{この章のまとめ}

\begin{itemize}
  \item コンピュータは、CPU・メモリ・ストレージ・GPU・入出力装置などの部品でできています。
  \item CPU は命令を実行する頭脳、メモリは作業中のデータを置く机、ストレージはデータを保存する本棚、GPU は大量の単純な計算を一斉にこなす専門家です。
  \item プログラムはストレージからメモリに読み込まれ、CPU が実行します。重い計算は GPU に任せることがあります。
  \item プログラムは、コンピュータに実行させる命令を順番に並べた指示書です。どんなプログラムも、順次・分岐・繰り返しの 3 つの組み合わせでできています。
  \item OS は、ハードウェアを管理し、アプリケーションが動く土台となるソフトウェアです。
  \item ロボットでは、Linux と ROS 2 を動かすメインのコンピュータと、モータを制御するマイコンを組み合わせて使うことがよくあります。
\end{itemize}
